\section{Maieutic Prompting: de la generación de explicaciones a la consistencia global}

La primera investigación parte de una ilusión común: el modelo puede dar una explicación, así que parece que «sabe» la respuesta. Sin embargo, el mismo modelo puede negar en la siguiente oración la razón que acaba de dar. Maieutic Prompting no confía directamente en la primera chain, sino que explora de forma recursiva explicaciones a favor y en contra, y después usa logic para elegir la asignación con mayor consistencia interna.

\subsection{Los Language Models solo son Sometimes Amazing}

El primer ejemplo le pide a GPT-3 que explique si, partiendo de la costa oeste de Estados Unidos y avanzando siempre hacia el oeste, se puede llegar finalmente a la costa este. El modelo da una explicación correcta basada en que la Tierra es redonda. Este acierto muestra que los modelos de lenguaje pueden combinar conocimiento implícito, pero no nos dice si el modelo es estable en preguntas cercanas. Al leer la figura, «responder bien» debe entenderse como un sample, no como una medición completa de los límites del conocimiento del modelo.

\lecturefigure{slide-03-language-models-sometimes-amazing.jpg}{Los modelos de lenguaje a veces generan explicaciones de sentido común coherentes y correctas}{Diapositiva V003 recuperada del video oficial; Stanford CS25 Lecture 16}

La siguiente diapositiva pone lado a lado preguntas true/false del mismo tipo. En ejemplos como «¿las mariposas tienen tres alas?», el modelo primero dice que tienen cuatro alas y después emite un juicio de verdadero o falso que contradice su propia explicación. Lo importante no es que cierto hecho de sentido común sea difícil, sino que entre las salidas falta un consistency contract: cada generación se comporta como una compleción local de texto, sin la obligación de mantener un belief state compartido.

\lecturefigure{slide-04-language-models-reasoning-failures.jpg}{Las preguntas cercanas exponen la contradicción entre la explanation y la final answer}{Diapositiva V004 recuperada del video oficial; Stanford CS25 Lecture 16}

\begin{warningbox}{Una explicación no equivale automáticamente a Faithful Reasoning}
La explanation que da un modelo de lenguaje puede ser solo texto generado junto con la respuesta. Aunque la explicación sea gramaticalmente correcta, puede contradecir otras explicaciones, citar hechos erróneos o invertirse cuando se reformula el prompt. Lo que Maieutic mejora en primer lugar es la consistencia lógica interna, no la verificación completa de la verdad sobre el mundo.
\end{warningbox}

\teachervoice{El ponente compara los modelos de lenguaje con «lemons» y luego dice que, si solo se eligen los ejemplos exitosos, parecen «cherries». Esta broma nos recuerda que la demo selection es en sí misma un evaluation bias. La verdadera pregunta no es si se puede encontrar una respuesta vistosa, sino si los errores tienen una estructura predecible.}

\subsection{Construcción recursiva del Explanation Tree}

El método Maieutic pide al modelo que genere explicaciones tanto para la proposición original $q$ como para su negación $\neg q$, y sigue preguntando por qué cada explicación se cumple o no. Lo que se obtiene no es una sola chain-of-thought, sino un argument graph local. Al leer la figura, primero hay que ubicar el nodo raíz y después comparar las ramas a favor, en contra y en conflicto mutuo; a mayor profundidad de búsqueda, hay más candidatos, pero también crecen el cómputo y el ruido.

\lecturefigure{slide-05-maieutic-logic-tree-example.jpg}{Maieutic tree: despliegue simultáneo de la proposición, su negación y las explicaciones recursivas}{Diapositiva V005 recuperada del video oficial; Jung et al., 2022}

La siguiente slide muestra la abductive generation: dada una conclusión, el modelo propone razones que podrían hacerla verdadera o falsa. Este paso aprovecha la amplitud generativa del LM y no exige que el modelo se comprometa de inmediato con una única respuesta. El sistema primero despliega el espacio del lenguaje en candidate explanation nodes y, en una etapa posterior, determina qué combinaciones pueden sostenerse a la vez.

\lecturefigure{slide-06-maieutic-abductive-generation.jpg}{Abductive generation: generación de explicaciones verificables para la proposición y su negación}{Diapositiva V006 recuperada del video oficial; Jung et al., 2022}

Sea $v_i$ un nodo de proposición y la asignación booleana $z_i\in\{0,1\}$ indique si el sistema la acepta al final. Si la explanation $e_j$ respalda la claim $c_j$, puede escribirse como una restricción de implicación
\begin{equation}
z_{e_j}\Rightarrow z_{c_j},
\end{equation}
mientras que la proposición y su negación explícita deben cumplir
\begin{equation}
z_q+z_{\neg q}=1.
\end{equation}
Aquí $z_q$ no es una probabilidad, sino una elección lógica tras la resolución global; la probabilidad solo se usa para ponderar las restricciones.

\begin{lstlisting}[caption={Pseudocódigo de explicación recursiva y consistencia global en Maieutic Prompting}]
def maieutic_reason(question, depth, language_model):
    graph = expand_both(question, negate(question), depth, language_model)
    clauses = []
    for explanation, claim, confidence in graph.edges:
        clauses.append(weighted_implication(explanation, claim, confidence))
    clauses.append(exactly_one(question, negate(question)))
    assignment = weighted_maxsat(clauses)
    return assignment[question], graph, assignment
\end{lstlisting}

\subsection{De la Local Confidence a Weighted MaxSAT}

El árbol recursivo aún puede contener ramas en conflicto mutuo, por lo que se necesita un global inference step. Weighted MaxSAT (máxima satisfacibilidad ponderada) busca, dentro de un conjunto de clauses booleanas, la asignación que maximiza el peso total satisfecho. No pretende que todas las oraciones se cumplan, sino que, cuando el conflicto es inevitable, conserva el conjunto más confiable y más compatible.

\lecturefigure{slide-07-maieutic-inference-objective.jpg}{Maieutic inference: conversión del explanation graph en una optimización lógica ponderada}{Diapositiva V007 recuperada del video oficial; Jung et al., 2022}

Si la $j$-ésima clause es $C_j(z)$ y su peso de confianza es $w_j$, el objetivo de resolución puede escribirse como
\begin{equation}
z^*=\arg\max_{z\in\{0,1\}^{|V|}}
\sum_j w_j\,\mathbf{1}\!\left[C_j(z)=\mathrm{true}\right].
\end{equation}
Una práctica común convierte la confianza del modelo $p_j$ en log-odds, $w_j=\log\frac{p_j}{1-p_j}$. $V$ es el conjunto de nodos del explanation graph y $\mathbf{1}[\cdot]$ es la función indicadora. La optimización garantiza consistencia relativa dadas las clauses; si todos los candidatos se basan en un sentido común erróneo, MaxSAT aún puede elegir mal de manera estable.

\begin{knowledgebox}{De qué se encarga lo Neural y de qué lo Symbolic}
El LM se encarga de la generación abierta y cubre las explicaciones que el modelo podría conocer; el symbolic solver se encarga de mantener las restricciones entre ramas. El primero resuelve el recall; el segundo mejora la consistency. Si solo se usa el solver, no hay conocimiento candidato; si solo se usa el LM, no hay consistencia global.
\end{knowledgebox}

\subsection{Resultados: al leer la gráfica de barras hay que fijarse en el Dataset y también en la Task}

La slide de resultados compara benchmarks true/false como CSQA 2.0, CREAK y Com2Sense. Maieutic supera a direct prompting y a varios baselines de self-consistency en varios conjuntos de datos, lo que indica que la explicación recursiva combinada con restricciones globales puede convertir el latent knowledge en respuestas de forma más estable. Al leer la figura, primero hay que comparar las diferencias entre métodos dentro de un mismo conjunto de datos y después observar el absolute level de cada conjunto; el muestreo y la difficulty de distintos benchmarks no permiten una comparación horizontal directa.

\lecturefigure{slide-08-maieutic-results.jpg}{Resultados de Maieutic Prompting en varios conjuntos de datos de juicios de verdadero o falso de sentido común}{Diapositiva V008 recuperada del video oficial; Jung et al., 2022}

\begin{warningbox}{Lo que esta figura no demuestra}
No demuestra que las explicaciones generadas sean faithful, ni que el modelo tenga un modelo completo del mundo; solo demuestra que, con estas tareas y esta configuración del prompt, buscar más explicaciones y añadir consistencia lógica mejora el juicio final. Para el despliegue también hay que considerar la latency, el token cost, el solver cost y la adversarial robustness.
\end{warningbox}

\subsection{Resumen de la sección}

Maieutic Prompting transforma una respuesta única en candidate generation más global consistency. Conserva la capacidad generativa del LM y, a la vez, reconoce que las salidas locales pueden entrar en conflicto. La conclusión didáctica más importante es esta: la estructura lógica puede mejorar la forma en que el modelo usa el conocimiento que ya tiene, pero no puede suplir de la nada el conocimiento del mundo que falta o que es erróneo.
